\section{实例分割}

\subsection{Mask R-CNN}

实例分割在目标检测的基础上，进一步为每个检测到的物体生成\textbf{像素级掩码}。\textbf{Mask R-CNN}（ICCV 2017）在 Faster R-CNN 的基础上增加了一个掩码预测分支：

\begin{figure}[H]
\centering
\includegraphics[width=0.95\textwidth]{slides-images/slide-054.jpg}
\caption{Mask R-CNN：在 Faster R-CNN 的基础上增加掩码预测分支，同时输出类别、边界框和实例掩码\protect\footnotemark}
\end{figure}
\footnotetext{Slides 第55页。}

\begin{enumerate}
    \item 对整张图运行 CNN + RPN 获得区域提议
    \item 对每个区域提取特征（\textbf{RoI Align}，使用双线性插值避免量化误差）
    \item 三个并行输出头：
    \begin{itemize}
        \item \textbf{分类头}：输出类别概率
        \item \textbf{回归头}：精调边界框坐标
        \item \textbf{掩码头}：使用 FCN 输出 $C \times M \times M$ 的二值掩码（$C$ 个类别，$M \times M$ 分辨率）
    \end{itemize}
\end{enumerate}

\begin{figure}[H]
\centering
\includegraphics[width=0.95\textwidth]{slides-images/slide-055.jpg}
\caption{Mask R-CNN 的检测效果：同时输出类别、边界框和精确的实例掩码\protect\footnotemark}
\end{figure}
\footnotetext{Slides 第56页。}

\begin{importantbox}{Mask R-CNN 的掩码预测机制}
掩码头为\textbf{每个类别}独立预测一个 $M \times M$ 的二值掩码（通常 $M=28$），最终根据分类头的预测选取对应类别的掩码作为输出。这种``分类与掩码解耦''的设计比``直接预测多类掩码''效果更好，因为它避免了类别间的竞争。掩码损失使用逐像素的二值交叉熵（Binary Cross-Entropy）。
\end{importantbox}

\begin{knowledgebox}{从 RoI Pooling 到 RoI Align 的关键改进}
原始 RoI Pooling 需要将浮点坐标量化为整数格点，这引入了最多半个格点的对齐误差。对于分类任务，这种误差可以容忍；但对于掩码预测，\textbf{哪怕一两个像素的偏差都会导致边界不准}。RoI Align 使用双线性插值替代硬量化，在非整数坐标处精确采样，使掩码精度显著提升。这一改进是 Mask R-CNN 成功的关键因素之一。
\end{knowledgebox}

\subsection{Panoptic Segmentation}

讲者还简要提到了\textbf{全景分割（Panoptic Segmentation）}——它统一了语义分割和实例分割：对``可数物体''（人、车、狗）做实例分割，对``不可数区域''（天空、草地、道路）做语义分割，最终输出一张完整的场景理解图。

\begin{table}[H]
\centering
\begin{tabular}{p{3cm}p{3.5cm}p{3.5cm}p{3cm}}
\toprule
\textbf{任务} & \textbf{输出} & \textbf{区分实例} & \textbf{代表方法} \\
\midrule
语义分割 & 每像素类别 & 否 & FCN, U-Net \\
实例分割 & 每实例掩码+类别 & 是（仅物体） & Mask R-CNN \\
全景分割 & 统一场景理解 & 是+语义背景 & Panoptic FPN \\
\bottomrule
\end{tabular}
\caption{三种分割任务对比}
\end{table}

\subsection{本章小结}

实例分割通过在目标检测基础上添加掩码预测，实现了像素级的实例区分。Mask R-CNN 凭借三头并行设计（分类 + 回归 + 掩码）和 RoI Align 的精确采样成为标杆方法。全景分割进一步统一了``物体''和``背景区域''的理解。

%% ========== Part 5: 可视化与理解 ==========

